% Figure: the Frank-Kamenetskii self-heating balance. The heat a
% slab loses by conduction (a line in the central-temperature excess)
% against the heat an Arrhenius reaction generates (an exponential),
% for three values of the generation parameter delta. Below the
% critical delta two steady states exist (a stable low one and an
% unstable high one); above it the generation curve clears the loss
% line everywhere and no steady state exists: thermal runaway.
\begin{figure}[htbp]
\centering
\begin{tikzpicture}
\begin{axis}[
  width=12cm, height=6.5cm,
  xlabel={dimensionless centre excess $\theta$},
  ylabel={dimensionless heat rate},
  xmin=0, xmax=6, ymin=0, ymax=6,
  grid=both,
  major grid style={engblack!20}, minor grid style={engblack!10},
  legend pos=north west, legend cell align=left,
  font=\small, samples=160, domain=0:6,
]
% Conductive loss: proportional to theta (linearised), slope 1.
\addplot[engblack, thick] { x };
\addlegendentry{conductive loss $\propto \theta$}
% Generation: delta * exp(theta) (Frank-Kamenetskii exponential).
% Sub-critical delta = 0.2: two crossings.
\addplot[engblue, very thick] { 0.2*exp(x) };
\addlegendentry{generation, $\delta = 0.2$ (sub-critical)}
% Critical delta_c = 1/e = 0.368: tangent to the loss line.
\addplot[engorange, very thick, dashed] { 0.3679*exp(x) };
\addlegendentry{generation, $\delta_c = 1/e$ (tangent)}
% Super-critical delta = 0.6: no crossing, runaway.
\addplot[engred, very thick, dotted] { 0.6*exp(x) };
\addlegendentry{generation, $\delta = 0.6$ (runaway)}
\end{axis}
\end{tikzpicture}
\caption[Thermal-runaway balance]{Self-heating balance for a
reacting slab (the Frank-Kamenetskii model). The straight line is
the heat conduction carries away, rising with the dimensionless
centre excess $\theta$; the curves are the heat an Arrhenius reaction
generates, $\delta\, e^{\theta}$, for three values of the generation
parameter $\delta$. At the sub-critical $\delta = 0.2$ the curve
crosses the loss line twice: the lower crossing is a stable steady
state, the upper an unstable one. At the critical $\delta_c = 1/e$
the curve is tangent to the line, the two steady states merge, and
any further increase leaves the generation above the loss for every
$\theta$, so no steady state exists and the temperature climbs
without bound. Crossing $\delta_c$ is the tipping point of thermal
runaway in a battery cell, a compost pile, or a stored reactive
solid, and the Poisson self-heating problem locates it before the
cell is built.}
\label{fig:vol02:ch12:runaway}
\end{figure}
